% !TEX program = xelatex
\IfFileExists{main.tex}{\documentclass[main.tex]{subfiles}}{
  \IfFileExists{../../main.tex}{\documentclass[../../main.tex]{subfiles}}{
    \IfFileExists{../main.tex}{\documentclass[../main.tex]{subfiles}}{
      \documentclass[Book\space NL/main.tex]{subfiles}
    }
  }
}
\begin{document}
\ifSubfilesClassLoaded{
  \ifcsname appendixchapterparent\endcsname
  \else
    \standaloneappendixchaptersetup{A}{2}
  \fi
}{}
\chapter{Membrane Transport and Cellular Excitability}
\margintoc

\section{Transport across cell membranes}

Cell membranes regulate exchange between the intracellular and extracellular
compartments. Small nonpolar molecules can diffuse through the lipid bilayer,
while ions and many polar solutes require channels or carrier proteins. Active
transport uses energy directly or indirectly to move substances against their
electrochemical gradients.

\section{Osmosis and cell volume}

Water moves across permeable membranes in response to differences in effective
solute concentration.\index{osmosis} Changes in extracellular tonicity can therefore alter
cell volume, while membrane transporters help maintain the composition of the
intracellular environment.

\section{Resting membrane potential}

The resting membrane potential\index{resting membrane potential} reflects ion concentration gradients and the
membrane's selective permeability. Potassium conductance is a major determinant
in many cells, while transporters maintain the underlying ion gradients.

\section{Action potentials}

In excitable cells, depolarization to threshold activates voltage-gated ion
channels and produces a rapid, regenerative change in membrane potential.\index{action potential}
Repolarization and refractory periods shape the timing and propagation of
electrical signals.
\end{document}